\subsection{扩张直线}

我们现在要通过给 R 添加两个新元素来定义一个拓扑空间 \(\overline{R}\) ，使得 R 到 R 的有界开区间 I 上的每个同胚都能延拓为 \(\overline{R}\) 到与 I 有相同端点的闭区间上的同胚。

设 \(\overline{\mathbf{R}}\) 是由（集合论， R, § 4 第 5 目）向 \(\mathbf{R}\) 添加两个新元素得到的集，这两个新元素记作 \(-\infty\) 与 \(+\infty\) 。我们通过把 \(\mathbf{R}\) 上的序延拓到 \(\overline{\mathbf{R}}\) ，方式是令 \(-\infty < a\) 与 \(a < +\infty\) （对所有 \(a \in \mathbf{R}\) ）以及 \(-\infty < +\infty\) ；显然我们如此得到一个全序集，其序在 \(\mathbf{R}\) 上诱导出实直线的序。其次，考察 \(\overline{\mathbf{R}}\) 上由 \(\overline{\mathbf{R}}\) 的开区间全体所生成的拓扑。由于在 \(\mathbf{R}\) 上 \(\overline{\mathbf{R}}\) 的开区间的迹是 \(\mathbf{R}\) 的开区间，这个拓扑在 \(\mathbf{R}\) 上诱导出实直线的拓扑。

定义 1. 赋有上述序结构与拓扑的集合 \(\overline{\mathbf{R}}\) 称为扩充实直线。

在使用扩张直线 \(\overline{\mathbf{R}}\) 时，出于用语习惯，把它的点称为实数往往是方便的；这时 \(\mathbf{R}\) 的点称为有限实数。我们在本节及本章随后三节中采用这一约定；今后凡采用这一约定时，我们都会明确指出它适用于正文的哪一部分。

若 \(a\) 是有限实数，则区间 \([a, +\infty[ \text{ 与 } ] - \infty, a]\) （相应地 \(]a, +\infty[ \text{ 与 } ] - \infty, a[\) ）（ \(\overline{\mathbf{R}}\) 中的）含于 \(\mathbf{R}\) ，且与 \(\mathbf{R}\) 中迄今记作 \([a, \rightarrow[ \text{ 与 } ] \leftarrow, a]\) （相应地 \(]a, \rightarrow[ \text{ 与 } ] \leftarrow, a[\) ）的区间一致；这个新记号用得远为频繁。再者，\(\mathbf{R}\) 与区间 \(] - \infty, +\infty[ \text{ 位于 } \overline{\mathbf{R}}, \text{ 且有时也如此记之。}\)

命题 2. 每个同胚 \(f\) （ \(\mathbf{R}\) 到区间 \(]a, b[\) 上的）都可以延拓为同胚 \(\overline{f}\) （ \(\overline{\mathbf{R}}\) 到 \([a, b]\) 上的）。若 \(f\) 是递增函数，则 \(\overline{f}\) 是 \(\overline{\mathbf{R}}\) 到 \([a, b]\) 上的一个序同构。

设 \(f\) 是递增同胚。若把 \(f\) 延拓到 \(\mathbf{R}\) ，方式是令 \(\overline{f}(-\infty) = a\) 与 \(\overline{f}(+\infty) = b\) ，则显然 \(\overline{f}\) 是 \(\overline{\mathbf{R}}\) 到 \([a, b]\) 上的严格递增映射（因而是双射）。于是 \(\overline{f}\) 把 \(\overline{\mathbf{R}}\) 的每个开区间映到关于 \([a, b]\) 为开的一个区间，因而由 § 1 第 4 目的定义 1 与命题 5，它是 \(\overline{\mathbf{R}}\) 到 \([a, b]\) 上的同胚。

若 \(f\) 递减，则对递增同胚 \(x \to -f(x)\) （ \(\mathbf{R}\) 到 \(] - b, - a[\) 上的）应用已证的结果。

因此，§ 2 中得到的区间 \([a, b]\) 的所有只涉及区间的序结构与拓扑的性质都可以搬到 \(\overline{R}\) 上；于是有下列命题：

命题 3. 扩充实直线是紧的。

因此（第二章 § 4 第 1 目，定理 1）在 \(\overline{\mathbf{R}}\) 上存在唯一一个与其拓扑相容的一致结构；这个一致结构同构于在 \([a, b]\) 上由 \(\mathbf{R}\) 的加法一致结构诱导的一致结构。但应当注意，在 \(\mathbf{R}\) 上由 \(\overline{\mathbf{R}}\) 的一致结构诱导的一致结构不是 \(\mathbf{R}\) 的加法一致结构（尽管它与实直线的拓扑相容）；因为 \(\mathbf{R}\) 关于其加法一致结构是完备空间，却不是 \(\overline{\mathbf{R}}\) 的完备子空间，因为它在 \(\overline{\mathbf{R}}\) 中不闭。

命题 4. \(\overline{\mathbf{R}}\) 的每个非空子集都有上确界与下确界。

R 的非空子集 A 的上确界（相应地下确界）记作 sup A（相应地 inf A）。显然

\[
\inf A \leqslant \sup A.\tag{1}
\]

若 \(\mathbf{A} \subset \mathbf{B}\) ，则 \(\sup \mathbf{A} \leqslant \sup \mathbf{B}\) 且 \(\inf \mathbf{A} \geqslant \inf \mathbf{B}\) （集合论，第三章 § 1 第 9 目，命题 4）。

命题 5. \(\overline{R}\) 的子集 A 连通当且仅当 A 是区间。

推论. 扩充实直线是连通的、局部连通的空间。

命题 6. 映射 \(f\) （区间 \(I\) （ \(\overline{R}\) 的）到 \(\overline{R}\) 中的）是 \(I\) 到 \(f(I)\) 上的同胚，当且仅当 \(f\) 在 \(I\) 上严格单调且连续；此时 \(f(I)\) 是 \(\overline{R}\) 的一个区间。

最后，函数 \(\sup (x,y)\) 与 \(\inf (x,y)\) 在 \(\overline{\mathbf{R}}\times \overline{\mathbf{R}}\) 上连续

